% Appendix companion to p1-main.tex; insert after the existing \appendix.
% Requires amsmath, amssymb, booktabs, and tabularx.
% This file intentionally does not issue \appendix itself.

\section{Additional Details of the Dependency-Driven Pipeline}
\label{app:p1}
This appendix supplements Section~\ref{sec:p1} with effect recognition,
dependence classification, and runtime realization details.

\subsection{Effect Extraction and Temporal Classification}
\label{app:p1-effects}
\paragraph{Loop identification and provenance.}
Special graph-neighbor loops are identified by their iterator structure.
Iterator-less loops are also classified as graph loops when an unseeded
provenance analysis of their stores finds a graph-data-derived GEP
subscript used by a store. Provenance propagates through SSA values to a
fixed point over the region domain
\[
\mathrm{Bottom}\sqsubseteq r\sqsubseteq\mathrm{Top},
\qquad r\in\{U,V,D,G\}.
\]
The four intermediate tags denote source, destination, data-derived, and
global regions. The region domain is an upper-bounded join-semilattice.

\paragraph{Effect atoms and structure.}
An effect atom is a tuple
\[
a=(\kappa,B,r,\omega,\tau,s),
\]
where $B$ is a canonical memory base,
$r\in\{\mathrm{Bottom},U,V,D,G,\mathrm{Top}\}$, and
$s\in\{\mathrm{Once},\mathrm{PerSource},\mathrm{PerPair}\}$.
The operation kind is
\[
\kappa\in\{R,W,U_f,U_{\oplus},\mathrm{Claim},\mathrm{Activate}\}.
\]
The optional $\omega$ identifies the operator of a recognized algebraic
update $U_{\oplus}$; $U_f$ denotes another old-value-dependent update.
A claim additionally carries its guard and transition. The annotation
$\tau$ records derived temporal information. Read analysis also records
whether the value source is live or a snapshot. Let $\Sigma$ denote the
finite multiset of primitive effect occurrences in a candidate loop.
Table~\ref{tab:effect-recognition} summarizes recognition.

The compiler groups occurrences into the preamble, neighbor body, and
epilogue of Equation~\ref{eq:effect-expr}. The empty effect $\epsilon$ is
the identity of both sequential composition $\triangleright$ and
structural unordered composition $\parallel$. For a finite domain
$H=(i_1,\ldots,i_n)$ in enumeration order,
\[
\mathrm{Seq}_{i\in H}X(i)
=X(i_1)\triangleright\cdots\triangleright X(i_n),
\]
with an empty domain yielding $\epsilon$. The domain size need not be
known at compile time. The binder path and phase identify each atom's
occurrence scope and position in the loop nest; these facts are retained
when constructing access footprints.

\begin{table*}[t]
\small
\centering
\caption{Recognition of effect atoms from SSA instructions and provenance.
Recognition identifies an operation pattern; it does not establish that
the operation can be reordered or parallelized.}
\label{tab:effect-recognition}
\begin{tabularx}{\textwidth}{lXX}
\toprule
\textbf{Kind} & \textbf{Recognition rule} & \textbf{Provenance and metadata} \\
\midrule
$R$
& A load from an address derived from $\operatorname{GEP}(A,i)$.
& The canonical base is $A$ and the region is $\pi(i)$. \\
$W$
& A store to an address derived from $\operatorname{GEP}(A,i)$.
& The canonical base is $A$ and the region is $\pi(i)$. \\
$U_{\oplus}$
& A store to $A[i]$ combines the previous value loaded from the same
location with another value using a recognized operator, such as
$+$, $-$, $\times$, bitwise AND, OR, XOR, or min/max. Guarded min/max
forms are also recognized.
& The load and store must refer to the same location. The index supplies
$r=\pi(i)$ and the recognized operator supplies $\omega$. \\
$U_f$
& A store depends on a load from the same location but does not match a
recognized algebraic-update pattern.
& The indexed location supplies the region $r$. \\
$\mathrm{Activate}$
& A next-frontier insertion. Set-based
\texttt{roaring\_bitmap\_add} is recognized directly; an array append
of the form $A[s{+}{+}]=v$ produces an activation atom.
& Explicit activation uses $\pi(v)$; the recognized array-append form
is assigned region $V$. \\
$\mathrm{Claim}$
& A conditional store of the form
\texttt{if (A[i] == expected) A[i] = desired}, with distinct expected
and desired values, is recognized as a first-wins claim pattern.
& The index supplies $r=\pi(i)$; guard operands and the state transition
are retained. \\
\bottomrule
\end{tabularx}
\end{table*}

\paragraph{Access conflicts and read summaries.}
For each read, the compiler compares mutations of the same base using
their regions, the mutation's phase, and the traversal domain. It
aggregates the comparisons into a four-valued temporal read summary,
using the highest-ranked relation. This summary is distinct from the
finer ordered RAW, WAR, and WAW relations used in interference analysis.

Access footprints also include loop-scoped runtime resources. Two
footprints on the same base or shared resource are potential conflict
candidates when at least one writes: the modes distinguish RAW, WAR,
and WAW. Read--read pairs do not conflict. A shared base alone does not
establish that two accesses reach the same element.

Phase order applies within a source iteration. Within a phase, the effect
multiset does not retain instruction order, so equal phase labels do not
establish separation or ordering. The internal read-before-write order
of a read--modify--write operation is retained explicitly.

\paragraph{Structural temporal classification.}
For a RAW or WAR candidate, let $r_R$ and $r_W$ be the read and write
regions and $p_W$ the write's phase. The following cases are applied in
order:
\begin{enumerate}
\item If $r_R=r_W$, classify the pair as \emph{SameRoundRead}.
\item For opposite endpoint regions $U$ and $V$, a write outside the
neighbor body gives \emph{SameRoundRead}. A neighbor-body write gives
\emph{PreviousRoundRead} under frontier traversal and \emph{Carried}
under all-vertex traversal.
\item For the remaining combinations, different phases give
\emph{SameRoundOrdered}; otherwise, this structural classification
infers no additional temporal relation.
\end{enumerate}
These classifications inform subsequent analysis rather than independently
authorizing a rewrite. In particular, a cross-endpoint relation depends on
the write's phase even when the read is in the preamble: per-source phase
order does not separate that read from other sources' neighbor updates.

WAW candidates have no read endpoint and are treated separately. When
the traversal cannot repeat across rounds, different-phase writes receive
\emph{SameRoundOrdered} and same-phase writes receive
\emph{SameRoundWrite}. When repeated rounds are possible, the cross-round
relation remains unknown until round analysis establishes the necessary
separation or ordering. This includes distinct dynamic occurrences of
the same write site.

\subsection{Dependence Relations and the Correctness Obligation}
\label{app:p1-independence}
\paragraph{Spatial and temporal footprints.}
Let $\mathcal L$ be the logical memory locations, including modeled
runtime resources. A spatial work instance has footprint
$S_i=(R_i,W_i)$, where $R_i,W_i\subseteq\mathcal L$.
A temporal execution unit, such as a round or phase component, similarly
has footprint $T_i=(R_i^T,W_i^T)$. These differ in their execution-unit
granularity, while using the same memory-conflict criterion.

For ordered units $x_i\prec x_j$, define
\[
\begin{split}
\mathsf{Conflict}(x_i,x_j)\iff{}&
W_i\cap R_j\ne\varnothing\ \lor\\
&R_i\cap W_j\ne\varnothing\ \lor\\
&W_i\cap W_j\ne\varnothing.
\end{split}
\]
The memory-dependence relation contains $(x_i,x_j)$ exactly when
$x_i\prec x_j$ and this predicate holds; temporal units use their
corresponding temporal footprints. Compiler approximations must
conservatively retain possible conflicts. Control-flow and lifecycle
requirements must also be preserved by the realization.

For a proposed non-trivial partition $\mathcal P$ of the units,
Equation~\ref{eq:p1-independence} states that no remaining dependence
crosses parts. Equivalently, an independent partition exists when the
remaining interference graph has more than one connected component
after ignoring edge direction. Order within each part remains subject
to the original dependences. A partition of temporal units and a
partition of spatial units are separate proof candidates; success for
one does not remove constraints imposed by the other.

\paragraph{Graph representation.}
The conceptual spatial graph has a node for each dynamic pair-work
instance $S(u,v)$ and conflict edges between instances; occurrence scopes
must additionally account for source-scoped work. At the temporal level,
phase nodes represent the preamble, body, and epilogue, with RAW, WAR,
and WAW relations between them. When inter-round interference is uniform,
the repeated relation can be summarized by one round node with a
self-loop. This cyclic summary is an interference graph, not an
execution DAG. Concrete ordered dependence instances and the emitted
task graph retain their required execution order.

\paragraph{Transformation discharge.}
For a transformation $L'=\mathcal T(L)$ and a proposed partition
$\mathcal P$, establishing
\[
\mathsf{Dep}_A(L')\cap\mathsf{Cross}(\mathcal P)=\varnothing
\]
establishes independence for the analyzed units of $L'$ on axis $A$.
It does not establish that $L'$ is equivalent to $L$. That is a separate
obligation under the transformation's applicability conditions, including
any requirements on update algebra, observations, claims, or snapshots.

\paragraph{Conditional preservation argument.}
Suppose that (i) the footprints conservatively describe all relevant
observable effects and shared resources; (ii) the selected transformations
preserve the original loop's observable sequential behavior; (iii) the
resulting units satisfy the independence criterion for the proposed
partition; and (iv) the realization preserves each part's computation,
required internal order, occurrence counts, and remaining control and
lifecycle constraints. Under the modeled sequential memory semantics,
nonconflicting units in different parts commute: neither writes a
location read or written by the other. A parallel interleaving can
therefore be rearranged into an execution of the transformed sequential
loop by exchanging adjacent independent units while preserving required
orders. Assumption (ii) then relates its result to the original loop.
This is the argument the analysis and lowering must support; the
independence equation alone is not a transformation-soundness theorem.

\subsection{Runtime Realization}
\label{app:p1-runtime}
\paragraph{Plans, callbacks, and task graphs.}
The realization $\mathcal R(L,E_t,\mathcal A)=\Pi_t$ is defined when the
available runtime mechanisms can implement the computation and its
ordering requirements. The compiler outlines computations into callbacks
and emits the information needed to instantiate a graph
$G_{\mathrm{exec}}=(V_{\mathrm{exec}},E_{\mathrm{exec}})$.
Runtime traversal sizes and partition counts determine the concrete task
instances. A node identifies an executable computation and its state;
an edge $x\rightarrow y$ requires $x$ to finish before $y$ starts.

The CPU runtime submits tasks to a common DAG executor. Tasks become
ready after all predecessors finish, are assigned to available workers,
and release successors on completion. Dispatch joins before returning
to its continuation. An edgeless graph represents independent tasks.
Task nodes can contain many effect occurrences and are not copied
directly from the interference graph. Phase order in $E_t$ need not
produce a separate task or stage for every phase.

\paragraph{Spatial execution.}
The compiler outlines the neighbor body into a pair callback receiving
the source $u$, destination $v$, and references to captured data,
including introduced snapshots or private storage. Supported preamble
and epilogue computations become source-entry and source-completion
callbacks. Accesses are redirected as required by the justified
transformations.

The compiler constructs or reuses the CleanCut layout and initializes
the traversal domain, callback arguments, and required snapshot, claim,
frontier, or private state. Source ownership places all neighbor work
for a source in one partition; destination ownership places all pair
work targeting a destination in one partition. The spatial proof
determines the required grouping, and CleanCut materializes a supported
partitioning consistent with it.

In source-owned execution, a partition task performs source entry,
neighbor computations, and source completion in order, retaining source
and neighbor enumeration order within the partition. Distinct partitions
can execute concurrently when their conflicts have been accounted for.
Destination ownership may split a source's neighbors between partitions,
so per-source operations cannot simply be repeated in each partition.
Where supported, source-coverage passes perform source entry before
pair traversal and source completion afterward, including for sources
with no neighbor pairs.

Standalone CPU spatial execution uses one task node per partition,
without precedence edges between independent partition tasks. If
destination ownership requires source preparation, a preparation node
precedes every partition task. Round initialization and snapshot
publication occur before dispatch. After the graph joins, the runtime
performs any destination-owned source-completion pass, combines private
results, and completes frontier publication and round finalization.
Source-owned completion callbacks instead execute inside partition
tasks. Operation descriptors identify callbacks and their lifecycle
events; a descriptor is not itself a parallel task.

\paragraph{Temporal execution.}
Temporal callbacks preserve control flow and instruction order within
their components. An execution frame stores captured values and local
state. For source-phase execution, dependent phases remain in one
component, and each independent component becomes a callback. The
runtime dispatches active components for the current source and joins
before advancing to the next source, making source-local state available
to the continuation.

For round execution, the compiler outlines a complete round body
parameterized by the round index. Rounds execute sequentially within
each group, while independent groups may run concurrently. Inputs are
captured and analysis-designated private state receives local storage.
Execution joins before leaving the original round loop, including
publication of the final counter value. Source-phase graphs have one
node per active component; round graphs have one node per round group.
Both graphs are edgeless because internal dependences stay within the
callbacks and the temporal proof establishes independence between nodes.

Supported independent source-owned and destination-owned computations
may also be emitted as child executions joined before their continuation.
A supported ordered realization executes the children in sequence.

\paragraph{Combined round and spatial execution.}
The compiler creates callbacks for reservation, round setup, partition
work, round finalization, and release. Each active round has its own
execution frame and mutable auxiliary state; graph layout may be shared.
For each round $t$, setup $S_t$, partition batches $B_{t,j}$, and
finalization $F_t$ obey
\[
S_t\longrightarrow B_{t,j}\longrightarrow F_t
\qquad(1\le j\le m).
\]
Setup executes the original computation preceding traversal and prepares
the partition resources. Each batch traverses its assigned partitions.
Finalization waits for all batches, performs combination and publication,
and executes the original computation following traversal.

Multiple independent rounds share an execution DAG, allowing overlap
within and across rounds. The runtime materializes a bounded set of
rounds and batches partitions according to the execution budget.
Resources are reserved before dispatch and released after the
corresponding execution joins. This combination reuses the established
round and spatial proof requirements; source-phase parallelism retains
the separate component realization.

\paragraph{Transformation state and lifecycle.}
Where justified, snapshotting redirects designated reads to a frozen
copy published before traversal. A staged per-source claim executes once
for each relevant source and gates its pair computations. Set-valued
frontiers retain current-frontier membership and suppress duplicate
activations when constructing the next frontier.

Eligible reductions and privatized updates receive private storage,
initialization, and combination operations. Partition partials are
combined in a specified order after parallel work completes; source-local
accumulations may instead be finalized at source completion. Validity
depends on the update and observation conditions established by the
analysis, including a valid combine operation and the absence of
unsupported observations of uncombined state.

\paragraph{GPU execution.}
The GPU backend uses the same analysis results and additionally checks
device support for computations and captured state. Supported callbacks
are lowered into device functions and kernel entry points. Descriptors
identify execution frames, referenced buffers, and pointer relocations.
The task DAG includes host operations, data publication, and kernel
launches. Device tasks use separate CUDA streams; successors are released
only after their required predecessors complete. Independent kernels
may overlap subject to device scheduling. Valid device buffers may be
reused, with transfers and synchronization performed as required.

For spatial execution, setup and finalization remain on the host. Setup
prepares state before device publication, spatial kernels execute the
partition work, and results become available before host finalization.
Where supported, independent ownership groups are distributed across
GPU threads while lifecycle and access order within each group are
preserved. Host finalization performs the required combination and
frontier publication.

For temporal execution, each independent component becomes a kernel
task. Currently, a component executes sequentially in one GPU thread,
so temporal concurrency occurs between component kernels. Source-phase
execution joins before advancing to the next source, and round execution
joins before leaving the original round loop.

Combined GPU execution materializes a bounded window of independent
rounds with separate execution frames. All host setup tasks precede a
shared upload $U_{\mathrm{pub}}$, which precedes every spatial kernel
$K_{t,j}$ in the window. A shared download $D_{\mathrm{pub}}$ waits for
all kernels and precedes every host finalization task. Thus, for every
active round and batch,
\[
S_t\longrightarrow U_{\mathrm{pub}}
\longrightarrow K_{t,j}\longrightarrow D_{\mathrm{pub}}
\longrightarrow F_t.
\]
Kernels may overlap across rounds and partitions, while shared transfer
boundaries introduce synchronization across the window. Capability and
offload selection determine whether this realization is used.
Unsupported device computations or unavailable offload use a supported
CPU realization before the affected computation executes.
